\documentclass[14pt,a4paper]{extreport}

\usepackage{iftex}
\ifPDFTeX
  \usepackage[T2A]{fontenc}
  \usepackage[utf8]{inputenc}
  \usepackage[russian,english]{babel}
\else
  \usepackage{fontspec}
  \IfFontExistsTF{Times New Roman}
    {\setmainfont{Times New Roman}}
    {\setmainfont{TeX Gyre Termes}}
  \IfFontExistsTF{Arial}
    {\setsansfont{Arial}}
    {\setsansfont{TeX Gyre Heros}}
  \IfFontExistsTF{Consolas}
    {\setmonofont{Consolas}}
    {\setmonofont{DejaVu Sans Mono}}
  \usepackage{polyglossia}
  \setdefaultlanguage{russian}
  \setotherlanguage{english}
\fi

\usepackage{geometry}
\geometry{left=30mm,right=15mm,top=20mm,bottom=20mm}

\usepackage{setspace}
\onehalfspacing

\usepackage{indentfirst}
\setlength{\parindent}{1.25cm}
\setlength{\parskip}{0pt}

\usepackage{array}
\usepackage{booktabs}
\usepackage{longtable}
\usepackage{enumitem}
\usepackage{caption}
\usepackage{titlesec}
\usepackage{tocloft}
\usepackage{xurl}
\usepackage[hidelinks,unicode]{hyperref}

\titleformat{\chapter}[block]
  {\normalfont\bfseries\centering\large}
  {\thechapter}
  {1em}
  {}
\titleformat{name=\chapter,numberless}[block]
  {\normalfont\bfseries\centering\large}
  {}
  {0pt}
  {}
\titleformat{\section}[block]
  {\normalfont\bfseries\normalsize}
  {\thesection}
  {1em}
  {}
\titlespacing*{\chapter}{0pt}{0pt}{18pt}
\titlespacing*{\section}{0pt}{12pt}{6pt}
\titlespacing*{\subsection}{0pt}{12pt}{6pt}

\renewcommand{\contentsname}{СОДЕРЖАНИЕ}
\renewcommand{\bibname}{СПИСОК ИСПОЛЬЗОВАННЫХ ИСТОЧНИКОВ}
\renewcommand{\figurename}{Рисунок}
\renewcommand{\tablename}{Таблица}
\captionsetup{labelsep=endash,justification=centering}

\setlist[itemize]{leftmargin=1.25cm,itemsep=0pt,topsep=0pt}
\setlist[enumerate]{leftmargin=1.25cm,itemsep=0pt,topsep=0pt}

\newcommand{\StudentName}{<ФИО студента>}
\newcommand{\StudentGroup}{<номер группы>}
\newcommand{\SupervisorName}{<ФИО, должность руководителя>}
\newcommand{\FacultyName}{<факультет/институт>}
\newcommand{\ProgramName}{<образовательная программа>}
\newcommand{\WorkTitle}{Разработка системы автоматизированного ревью merge request на основе генерации и ранжирования комментариев}

\newcommand{\unnumberedchapter}[1]{%
  \chapter*{#1}%
  \addcontentsline{toc}{chapter}{#1}%
}

\begin{document}

\begin{titlepage}
  \thispagestyle{empty}
  \begin{center}
    Министерство науки и высшего образования Российской Федерации\\
    Федеральное государственное автономное образовательное учреждение\\
    высшего образования\\[0.5em]
    \textbf{НАЦИОНАЛЬНЫЙ ИССЛЕДОВАТЕЛЬСКИЙ УНИВЕРСИТЕТ ИТМО}\\[2em]

    \FacultyName\\
    \ProgramName
  \end{center}

  \vspace{3em}

  \begin{center}
    \textbf{ОТЧЕТ}\\[0.5em]
    по научно-исследовательской работе\\[2em]
    на тему:\\[0.5em]
    \textbf{\WorkTitle}\\[0.5em]
    \textbf{(проект MergeMind)}
  \end{center}

  \vfill

  \begin{flushright}
    Выполнил: студент группы \StudentGroup\\
    \StudentName\\[1.5em]
    Руководитель НИР:\\
    \SupervisorName
  \end{flushright}

  \vfill

  \begin{center}
    Санкт-Петербург\\
    2026
  \end{center}
\end{titlepage}

\setcounter{page}{2}

\unnumberedchapter{РЕФЕРАТ}

Отчет по научно-исследовательской работе содержит черновое описание проекта
MergeMind: локально запускаемой системы для автоматизированного ревью merge
request. Объектом исследования является процесс проверки изменений в
программном коде, предметом исследования --- методы генерации, отбора и
проверки полезности review-комментариев по контексту изменения.

Цель работы --- разработать и экспериментально проверить MVP-пайплайн, который
по diff и контексту репозитория формирует небольшое число конкретных,
обоснованных и практически применимых замечаний. Для достижения цели
рассматриваются задачи подготовки MR-centric данных, извлечения структурного
контекста, построения генератора и ранжировщика комментариев, подключения
локальных LLM-моделей, а также оценки качества через offline-метрики,
LLM-judge и downstream-проверку в SWE-CI.

В текущей версии реализованы адаптеры датасетов CodeReviewer, CodeReviewQA и
CoDocBench, модуль разбора unified diff, retrieval baseline, LLM-компоненты
generator--reranker--rewriter, экспериментальные режимы для GitHub PR и SWE-CI,
а также мониторинг и сохранение воспроизводимых артефактов. Первичные
эксперименты показывают, что проект пока не доказывает сокращение числа
итераций CI-loop, но дает положительные сигналы по groundedness, usefulness и
снижению финального test gap на отдельных SWE-CI smoke-задачах.

\textbf{Ключевые слова:} automated code review, merge request, large language
models, reranking, LLM judge, SWE-CI, GitHub pull request, diff analysis.

\clearpage
\tableofcontents
\clearpage

\unnumberedchapter{ПЕРЕЧЕНЬ СОКРАЩЕНИЙ И ОБОЗНАЧЕНИЙ}

\begin{longtable}{>{\raggedright\arraybackslash}p{0.22\textwidth}
                  >{\raggedright\arraybackslash}p{0.68\textwidth}}
  \toprule
  Сокращение & Расшифровка \\
  \midrule
  MR & Merge request, запрос на внесение изменений в кодовую базу. \\
  PR & Pull request, аналог MR в GitHub. \\
  LLM & Large language model, большая языковая модель. \\
  MVP & Minimum viable product, минимально работоспособная версия продукта. \\
  CI & Continuous integration, непрерывная интеграция. \\
  SWE-CI & Benchmark для оценки coding-agent в цикле непрерывной интеграции. \\
  LLM judge & Оценивание результатов с помощью языковой модели по заданной рубрике. \\
  \bottomrule
\end{longtable}

\chapter*{ВВЕДЕНИЕ}
\addcontentsline{toc}{chapter}{ВВЕДЕНИЕ}

Современная разработка программного обеспечения почти всегда включает
обязательное ревью изменений. Ревью позволяет обнаруживать дефекты, снижать
архитектурные риски, поддерживать единый стиль разработки и распространять
знание о кодовой базе внутри команды. Однако ручная проверка merge request
требует времени опытных разработчиков, а качество ревью зависит от загруженности
команды, полноты контекста и способности ревьюера быстро понять последствия
изменения.

Автоматизация code review развивалась от статических анализаторов и правил
линтинга к моделям, которые работают с diff, естественно-языковыми
комментариями и контекстом репозитория. При этом простая генерация большого
числа замечаний не решает задачу: шумные или непроверенные комментарии снижают
доверие к инструменту и увеличивают нагрузку на разработчика. Поэтому в проекте
MergeMind основной акцент сделан на отборе малого числа замечаний с высокой
ожидаемой практической пользой.

\textbf{Цель работы} --- разработать локально запускаемый пайплайн
автоматизированного ревью MR, который по контексту изменения генерирует и
отбирает 1--3 наиболее полезных комментария.

Для достижения цели поставлены следующие задачи:

\begin{enumerate}
  \item изучить существующие подходы к automated code review и релевантные
    датасеты;
  \item привести разные источники данных к единой MR-centric схеме;
  \item реализовать обработку контекста изменений: diff, файлы, hunks,
    добавленные и удаленные строки, changed identifiers;
  \item построить baseline-пайплайн генерации и ранжирования комментариев;
  \item подключить локальные LLM-компоненты через LM Studio / Qwen;
  \item разработать процедуру оценки качества комментариев;
  \item провести первичные offline, GitHub PR и SWE-CI эксперименты.
\end{enumerate}

\textbf{Объект исследования} --- процесс ревью изменений в программном коде.
\textbf{Предмет исследования} --- методы генерации, ранжирования и проверки
полезности review-комментариев на основе diff и контекста репозитория.

\textbf{Практическая значимость} работы состоит в возможности использовать
MergeMind как локальный ассистент ревьюера или как дополнительный слой проверки
coding-agent в CI-like workflow. Важным ограничением проекта является запрет на
использование скрытого target-fix в ревью: комментарии должны строиться только
на доступном diff, требованиях и видимом контексте.

\chapter{ОБЗОР ПРЕДМЕТНОЙ ОБЛАСТИ}

\section{Роль code review в разработке}

Code review используется для проверки корректности изменений до их попадания в
основную ветку проекта. В отличие от линтинга или unit-тестов, ревью учитывает
семантику задачи, совместимость с существующей архитектурой, потенциальные
регрессии, читаемость и сопровождаемость кода. Эти свойства делают ревью
трудно автоматизируемым: полезный комментарий должен быть не только похож на
комментарии из истории, но и быть привязанным к конкретному риску изменения.

Практическая ценность automated code review определяется не количеством
сгенерированных замечаний, а их применимостью. Если инструмент регулярно
предлагает несущественные, спорные или неоснованные комментарии, разработчики
начинают игнорировать его вывод. Следовательно, для MergeMind важны три
критерия: groundedness, usefulness и actionability.

\section{Существующие подходы к автоматизации ревью}

Ранние подходы к автоматизации ревью часто опирались на правила, статический
анализ и поиск известных паттернов ошибок. Такие инструменты хорошо обнаруживают
локальные нарушения, но плохо работают с проектно-специфичным контекстом и
естественно-языковыми объяснениями.

Модели семейства CodeReviewer рассматривают code review как отдельную область
предобучения и решают задачи quality estimation, comment generation и code
refinement. Это важно для MergeMind, поскольку проект также работает с diff и
review comments, но делает дополнительный акцент на ранжировании и отсечении
шума.

Другой класс работ рассматривает не только генерацию текста, но и способность
модели понять замечание ревьюера, связать его с изменением и предложить
исправление. В этом направлении полезен CodeReviewQA, так как он смещает фокус
с текстового сходства на reasoning вокруг review-комментария.

\section{Данные и контекст изменений}

Для MVP используются три источника данных:

\begin{itemize}
  \item CodeReviewer --- основной источник исторических diff и review comments;
  \item CodeReviewQA --- validation-side источник для проверки reasoning;
  \item CoDocBench --- вспомогательный источник для change alignment и анализа
    связей между изменениями кода и сопутствующими артефактами.
\end{itemize}

Единица обработки в MergeMind --- пример MR. Он включает идентификатор,
метаданные, unified diff, список измененных файлов, контекст репозитория,
исторические gold comments и опциональные сигналы результата. Такой формат
позволяет запускать один и тот же inference pipeline на offline-датасетах,
публичных GitHub PR и SWE-CI задачах.

\section{Оценка качества review-комментариев}

Текстовое сходство с human review полезно как базовый сигнал, но оно
недостаточно для оценки практической пользы. Реальный ревьюер мог не заметить
ошибку, оставить комментарий другой формулировкой или сфокусироваться на
локальном стиле, тогда как модель может указать на иной, но обоснованный риск.

Поэтому в MergeMind используются несколько уровней оценки:

\begin{itemize}
  \item deterministic similarity, hit@k и MRR для сопоставления с gold comments;
  \item LLM judge с рубриками groundedness, usefulness, valid alternative и
    gold alignment;
  \item ручной просмотр side-by-side отчетов;
  \item downstream-проверка в SWE-CI, где review-style комментарий может
    влиять на поведение coding-agent.
\end{itemize}

\chapter{ПОСТАНОВКА ЗАДАЧИ И АРХИТЕКТУРА MERGEMIND}

\section{Функциональные требования}

Система должна принимать на вход MR или PR, извлекать из него diff и
доступный контекст, генерировать candidate-комментарии, ранжировать их по
ожидаемой полезности и возвращать малый набор финальных замечаний. Для MVP
достаточно локального запуска из командной строки, без публикации комментариев
обратно в GitHub.

Функциональные требования:

\begin{enumerate}
  \item поддержка подготовки данных из нескольких источников;
  \item единая схема MRExample для offline и live-сценариев;
  \item обработка unified diff и извлечение структурных признаков;
  \item baseline без внешних API;
  \item LLM-пайплайн через локальный OpenAI-compatible endpoint;
  \item сохранение predictions, metrics и конфигурации запуска;
  \item безопасный режим GitHub PR demo без записи в репозиторий;
  \item интеграция с SWE-CI без использования target\_sha в review context.
\end{enumerate}

\section{Нефункциональные требования}

Для исследовательского MVP важны воспроизводимость, наблюдаемость и
контролируемая стоимость экспериментов. Каждый запуск должен сохранять
конфигурацию, прогнозы, метрики, логи и сведения о фактической модели.
Ошибки парсинга и fallback-режимы должны учитываться явно, так как silent
fallback может создать ложное впечатление успешного результата.

Отдельное требование --- локальность. Проект рассчитан на запуск baseline без
платных API, а LLM-эксперименты могут выполняться через LM Studio и локальные
модели Qwen. При использовании внешних провайдеров ключи хранятся только во
внепроектном окружении или в файле \texttt{.env}, который не должен попадать в
git.

\section{Общая архитектура}

Архитектура MergeMind строится как последовательный pipeline:

\begin{center}
\begin{tabular}{>{\raggedright\arraybackslash}p{0.24\textwidth}
                >{\raggedright\arraybackslash}p{0.62\textwidth}}
  \toprule
  Блок & Назначение \\
  \midrule
  Data adapters & Загрузка CodeReviewer, CodeReviewQA, CoDocBench и приведение к MRExample. \\
  Context processing & Разбор diff, hunks, changed files, identifiers и repository context. \\
  Generator & Формирование candidate review comments. \\
  Reranker & Отбор наиболее полезных и обоснованных комментариев. \\
  Rewriter & Сжатие и нормализация финального текста комментария. \\
  Validation & Расчет метрик, LLM judge, отчеты для ручного просмотра. \\
  Monitoring & Снимки состояния runs, dashboard и chronicle. \\
  \bottomrule
\end{tabular}
\end{center}

Такое разделение позволяет независимо сравнивать генераторы, ранжировщики и
rewrite-профили, не меняя формат данных и общий evaluation flow.

\section{Сценарии использования}

Первый сценарий --- offline evaluation на подготовленных train/validation/test
данных. Он нужен для быстрого сравнения baseline и LLM-профилей.

Второй сценарий --- GitHub PR demo. Система загружает публичный PR в read-only
режиме, строит MRExample, запускает inference pipeline и сохраняет отчет. Если
в PR уже есть human review comments, они используются как временный gold для
метрик сходства.

Третий сценарий --- SWE-CI validation. Здесь MergeMind используется как
дополнительный review layer для coding-agent. Важная часть постановки состоит в
том, что target commit не передается в prompt или review artifacts: система
работает только с текущим patch/diff и доступным контекстом.

В текущей реализации комментарий MergeMind передается не после завершения
SWE-CI, а внутри epoch: сначала architect формирует \texttt{requirement.xml},
затем programmer делает первичную правку, после этого MergeMind анализирует
diff между состоянием до и после programmer patch и формирует review comment.
Комментарий копируется в контейнер как \texttt{/app/mergemind\_review.md};
рядом передаются \texttt{/app/requirement.xml}, список разрешенных файлов для
revision и, в отдельных профилях, before-snapshot только тех файлов, которые
уже изменил programmer. После этого programmer выполняет revision pass, и
pytest запускается уже на revised-коде. Поэтому эксперимент проверяет именно
in-loop влияние комментария, а не post-hoc оценку готового решения.

\chapter{РЕАЛИЗАЦИЯ MVP}

\section{Структура репозитория}

Код проекта разделен на несколько модулей:

\begin{itemize}
  \item \texttt{src/data} --- загрузка и нормализация датасетов;
  \item \texttt{src/context} --- diff parsing и lightweight structural context;
  \item \texttt{src/models} --- retrieval baseline, reranker и LLM-компоненты;
  \item \texttt{src/validation} --- offline-метрики, LLM judge и SWE-CI слой;
  \item \texttt{src/inference} --- сборка полного inference flow;
  \item \texttt{src/monitoring} --- dashboard и monitoring agent;
  \item \texttt{scripts} --- CLI-команды для подготовки данных, экспериментов и
    отчетности.
\end{itemize}

Такая структура отделяет исследовательские эксперименты от ядра pipeline и
упрощает добавление новых режимов без переписывания общей схемы.

\section{Подготовка данных}

Адаптеры датасетов приводят разные форматы к единой структуре MRExample.
Для CodeReviewer используются реальные diff и review comments. Для
CodeReviewQA создается синтетический unified diff на основе старой и новой
версии фрагмента кода, а reference review сохраняется как gold comment. Для
CoDocBench учитываются связи между изменениями кода и документацией.

На этапе подготовки задаются лимиты длины diff, repository context и файлового
контекста. Это необходимо, чтобы inference оставался воспроизводимым и не
зависел от случайного переполнения контекстного окна LLM.

\section{Обработка diff и контекста}

Модуль \texttt{src/context/processing.py} разбирает unified diff, выделяет
измененные файлы, hunks, добавленные и удаленные строки. Для структурного
обогащения используется Tree-sitter как best-effort механизм. Если parser не
доступен или не справляется с конкретным языком, pipeline не падает, а
продолжает работу на уровне файла и hunk.

Результатом обработки является repository context: компактное текстовое
представление измененных файлов, функций и идентификаторов, которое можно
передать генератору и ранжировщику.

\section{Модельный pipeline}

Базовая схема имеет вид:

\begin{center}
\texttt{context -> generator -> reranker -> optional rewriter -> validation}
\end{center}

Retrieval baseline ищет похожие исторические примеры и предлагает комментарии
из ближайших MR. Logistic reranker и heuristic fallback оценивают признаки
конкретности, связи с diff, упоминания ошибок, тестов и потенциальных
регрессий.

LLM-профили расширяют baseline: генератор формирует candidate comments,
reranker оценивает их полезность, а rewriter делает итоговый комментарий
коротким и пригодным для ревью. Для SWE-CI добавлены специальные prompt-профили,
которые фокусируются не на стиле, а на correctness-only рисках и ожидаемом
repair impact.

\section{Валидация и мониторинг}

Offline-валидация считает similarity, hit@k, MRR, latency и runtime-метрики.
LLM judge используется для оценки gold alignment, valid alternative,
groundedness и usefulness. Для ручной проверки генерируются side-by-side
отчеты с diff, gold comments и predictions.

Monitoring agent сохраняет chronicle, dashboard, snapshot и presentation
artifacts. Это важно для НИР, так как экспериментальный проект развивается
итеративно, а промежуточные выводы нужно фиксировать вместе с конфигурациями и
ограничениями запусков.

\chapter{ЭКСПЕРИМЕНТАЛЬНАЯ ОЦЕНКА}

\section{Offline и GitHub PR эксперименты}

В offline-режиме проект сравнивает baseline и LLM-профили по подготовленным
данным. Однако для практической оценки автоматического ревью важны примеры из
реальных PR. Поэтому была подготовлена выборка публичных merged GitHub PR, для
которых MergeMind строил комментарии в read-only режиме.

Сравнение текущего \texttt{qwen35\_rewriter} с contract-профилем показало, что
contract-подход лучше фокусируется на grounded correctness risks.

\begin{center}
\begin{tabular}{lrrrrr}
  \toprule
  Pipeline & PR & Comments & Avg judge & Grounded & Useful \\
  \midrule
  \texttt{qwen35\_rewriter} & 5 & 12 & 0.650 & 0.730 & 0.660 \\
  \texttt{qwen35\_review\_contract} & 5 & 15 & 0.794 & 0.920 & 0.790 \\
  \bottomrule
\end{tabular}
\end{center}

При этом hit@k остался равен нулю для обеих конфигураций. Это не следует
интерпретировать как полный провал, потому что сгенерированные замечания могут
быть полезными альтернативами, не совпадающими лексически с human review.
Например, в реальных PR система находила риск зависания после удаления timeout
logic, ошибочный порядок вызовов ORM и потенциальную проблему rate limit при
параллельных запросах.

\section{SWE-CI эксперименты}

SWE-CI используется как downstream benchmark: проверяется, помогает ли
review-style комментарий coding-agent двигаться к меньшему test gap. На текущем
этапе строгий критерий сокращения числа итераций не достигнут: baseline и
assisted runs в smoke-задачах обычно использовали все разрешенные итерации.
Тем не менее появились положительные сигналы по финальному gap и предотвращению
регрессий.

\begin{center}
\begin{tabular}{p{0.34\textwidth}rrrrr}
  \toprule
  Run & EvoScore & Final gap & Best gap & Comments & Revisions \\
  \midrule
  Baseline \texttt{inline-snapshot} & 0.5045 & 12 & 12 & 0 & 0 \\
  Contract \texttt{inline-snapshot} & 0.9167 & 13 & 12 & 9 & 3 \\
  Triage \texttt{inline-snapshot} & 0.8333 & 12 & 12 & 2 & 2 \\
  Baseline \texttt{cle-b/httpdbg} & -0.5439 & 11 & 5 & 0 & 0 \\
  Triage-control \texttt{cle-b/httpdbg} & 0.0000 & 5 & 5 & 1 & 1 \\
  \texttt{qwen35\_caveman\_top1} & -0.3333 & 5 & 5 & 2 & 2 \\
  \bottomrule
\end{tabular}
\end{center}

На задаче \texttt{inline-snapshot} положительный сигнал виден по официальной
метрике SWE-CI: baseline без MergeMind получил EvoScore 0.5045, тогда как
contract- и triage-профили MergeMind получили 0.9167 и 0.8333 соответственно.
Это означает, что внутри заданного числа итераций assisted-варианты сильнее
приближались к целевому числу проходящих тестов.

На задаче \texttt{cle-b/httpdbg} положительного EvoScore пока не получено, но
assisted-профили улучшили final gap с 11 до 5. В этом кейсе важно различать
численный gap и конкретный набор падающих тестов, поэтому дополнительно
считались failed-test-set overlap, fixed failures и new failures. Для лучшего
caveman-профиля новых финальных падений не появилось, а 6 финальных падений
baseline отсутствовали в assisted final set. Строгий revision guard также дал
проверяемый safety-сигнал: он отклонил попытку изменить файл вне programmer
patch. Последующая before-snapshot версия добавила в контекст только содержимое
файлов до текущего programmer patch, не передавая target/oracle information.

\section{Интерпретация результатов}

Текущие эксперименты подтверждают, что задача автоматического ревью не сводится
к генерации похожего текста. В GitHub PR сценарии LLM judge показывает рост
groundedness и usefulness при переходе от общего rewriter к contract-профилю. В
SWE-CI сценарии отдельные assisted-конфигурации улучшают EvoScore или уменьшают
final gap, но пока не сокращают число итераций.

Главные ограничения текущей оценки:

\begin{itemize}
  \item малый размер SWE-CI smoke-выборки;
  \item зависимость от локальной модели и параметров LM Studio;
  \item неполное совпадение human review comments с полезными альтернативными
    замечаниями;
  \item высокая стоимость многоагентных SWE-CI запусков;
  \item необходимость ручной проверки части LLM judge выводов.
\end{itemize}

\chapter{ЗАКЛЮЧЕНИЕ}

В ходе НИР начата разработка MergeMind --- системы автоматизированного ревью
merge request, ориентированной на малое число полезных комментариев вместо
массовой генерации замечаний. Реализованы подготовка данных, обработка diff,
retrieval baseline, LLM-компоненты generator--reranker--rewriter, GitHub PR demo,
SWE-CI слой, monitoring agent и набор тестов.

Основной результат текущего этапа --- сформирована воспроизводимая
исследовательская инфраструктура, в которой можно сравнивать разные профили
автоматического ревью не только по текстовым метрикам, но и по downstream
поведению coding-agent. На публичных PR contract-профиль улучшил средние
оценки judge, groundedness и usefulness. В SWE-CI smoke-задачах отдельные
конфигурации улучшили final gap, однако строгий критерий сокращения числа
итераций пока не выполнен.

Дальнейшая работа должна быть направлена на расширение SWE-CI выборки,
калибровку LLM judge по ручной разметке, улучшение context retrieval,
добавление test-failure context как first-class сигнала и снижение стоимости
review-профилей. Отдельным направлением является подготовка полноценного отчета
с финальными таблицами экспериментов, приложениями и заполненными
административными данными титульного листа.

\begin{thebibliography}{99}
\addcontentsline{toc}{chapter}{СПИСОК ИСПОЛЬЗОВАННЫХ ИСТОЧНИКОВ}

\bibitem{gost732}
ГОСТ 7.32--2017. Система стандартов по информации, библиотечному и
издательскому делу. Отчет о научно-исследовательской работе. Структура и
правила оформления.

\bibitem{itmo-nir}
Университет ИТМО. Отчет о НИР: материалы для руководителя проекта.
Раздел сайта ``Исследования и разработки''. URL:
\url{https://science.itmo.ru/} (дата обращения: 25.05.2026).

\bibitem{itmo-method}
Университет ИТМО. Методические материалы по подготовке отчета по НИР
магистранта. URL: \url{https://books.ifmo.ru/file/pdf/3000.pdf}
(дата обращения: 25.05.2026).

\bibitem{codereviewer}
Li Z., Lu S., Guo D. et al. Automating Code Review Activities by
Large-Scale Pre-training. arXiv:2203.09095, 2022.
URL: \url{https://arxiv.org/abs/2203.09095}.

\bibitem{towards-acr}
Tufano M., Watson C., Bavota G., Poshyvanyk D. Towards Automating Code Review
Activities. arXiv:2101.02518, 2021.
URL: \url{https://arxiv.org/abs/2101.02518}.

\bibitem{codereviewqa}
CodeReviewQA: The Code Review Comprehension Assessment for Large Language
Models. Findings of ACL, 2025. URL:
\url{https://aclanthology.org/2025.findings-acl.476.pdf}.

\bibitem{codocbench}
CoDocBench: A Dataset for Code-Documentation Alignment in Software Maintenance.
arXiv:2502.00519, 2025. URL: \url{https://arxiv.org/abs/2502.00519}.

\bibitem{geval}
Liu Y., Iter D., Xu Y., Wang S., Xu R., Zhu C. G-Eval: NLG Evaluation using
GPT-4 with Better Human Alignment. arXiv:2303.16634, 2023.
URL: \url{https://arxiv.org/abs/2303.16634}.

\bibitem{sweci}
SWE-CI: Evaluating Agent Capabilities in Maintaining Codebases via Continuous
Integration. arXiv:2603.03823, 2026.
URL: \url{https://arxiv.org/abs/2603.03823}.

\bibitem{treesitter}
Tree-sitter: An incremental parsing system for programming tools.
URL: \url{https://github.com/tree-sitter/tree-sitter}
(дата обращения: 25.05.2026).

\end{thebibliography}

\appendix

\chapter{КОМАНДЫ ДЛЯ ВОСПРОИЗВЕДЕНИЯ ЭКСПЕРИМЕНТОВ}

Ниже приведены базовые команды, использованные как основа для запуска
локальных экспериментов.

\section*{Подготовка данных}

\begin{verbatim}
python scripts/download_datasets.py
python scripts/prepare_data.py
\end{verbatim}

\section*{Baseline без LLM}

\begin{verbatim}
python scripts/train_baseline.py
python scripts/evaluate.py
python scripts/demo_mr.py
\end{verbatim}

\section*{LLM-пайплайн}

\begin{verbatim}
python scripts/check_llm.py
python scripts/evaluate.py --pipeline qwen35_rewriter --profile smoke
\end{verbatim}

\section*{GitHub PR experiment}

\begin{verbatim}
python scripts/run_github_pr_experiment.py \
  --manifest configs/pr_experiments/curated_real_prs.json \
  --config configs/base.yaml
\end{verbatim}

\section*{SWE-CI assisted smoke}

\begin{verbatim}
python scripts/run_swe_ci.py \
  --swe-ci-repo-path ~/SWE-CI \
  --tasks-path configs/swe_ci_caveman_low_gap_tasks.jsonl \
  --output-dir artifacts/swe_ci_runs \
  --run-id caveman_grid_top1_cle_b_001 \
  --limit 1 \
  --max-iterations 3 \
  --mode mergemind_assisted \
  --agent-name direct_openai \
  --mergemind-pipeline qwen35_caveman_top1 \
  --mergemind-top-n 1
\end{verbatim}

\chapter{ПЛАН ДОРАБОТКИ ОТЧЕТА}

\begin{enumerate}
  \item Заполнить титульный лист: факультет, программу, ФИО, группу и
    руководителя.
  \item Уточнить требуемый тип отчета: семестровая НИР, аннотированный отчет
    или отчет по проекту.
  \item Добавить финальные таблицы по всем завершенным экспериментам.
  \item Проверить библиографию под требования конкретного подразделения ИТМО.
  \item После установки TeX Live или MiKTeX собрать PDF и исправить переносы,
    таблицы и нумерацию страниц.
\end{enumerate}

\end{document}
